% ─────────────────────────────────────────────────────────────────────────────
% Abstract — 8 required components for IEEE TIFS
% ─────────────────────────────────────────────────────────────────────────────

\begin{abstract}
% [1] Motivation / Problem
Backdoor attacks on large language models (LLMs) persist through safety
fine-tuning procedures, yet no principled metric exists to quantify this
persistence as a function of model scale.
% [2] Gap / What is missing
Existing attack success rate (ASR) measurements are scale-agnostic and provide
no theoretical guarantees connecting parameter count to backdoor robustness.
% [3] Proposed solution / Novelty
We propose the \emph{Cyber Resilience Scaling Coefficient} (CRSC), a closed-form
metric that combines normalized Frobenius drift of hidden-state representations
($\Delta_\text{hidden}$) and Shannon entropy change of the loss distribution
($\Delta H$) into a single score quantifying backdoor persistence after
supervised fine-tuning (SFT) or reinforcement learning from human feedback (RLHF).
% [4] Theoretical contribution
We prove two properties: (i) CRSC is monotonically non-decreasing in parameter
count $N$ for fixed fine-tuning conditions (Proposition 1), formalized via
random matrix theory; and (ii) CRSC converges to a concentration-of-measure
proxy for ASR with bounded estimation error (Proposition 2).
% [5] Method / Implementation
A three-agent evaluation pipeline (Detection, Analysis, Response) computes CRSC
across $N \in \{7\text{B}, 13\text{B}, 70\text{B}\}$, poison rates
$\rho \in \{0.001, 0.005, 0.01, 0.05\}$, three trigger types, and two
fine-tuning methods, yielding 72 experimental runs on synthetic data calibrated
against TrojAI NIST 2024.
% [6] Results
At threshold $\tau = 0.62$, CRSC achieves F1\,$=\,0.891$ (95\% CI: $[0.874, 0.908]$),
AUC-ROC\,$=\,0.924$, and Pearson $r\,=\,0.847$ ($p < 10^{-5}$) between CRSC and
$\log_{10}(N)$, confirming the scaling monotonicity.
% [7] Compliance
CRSC is mapped to MITRE ATLAS adversarial ML techniques and the NIST AI Risk
Management Framework, providing a compliance-ready risk classification at three
severity levels.
% [8] Availability
Code and data are released under MIT license at
\url{https://github.com/ufxa/Cyber-Resilience-Scaling-Coefficient}.
\end{abstract}

\begin{IEEEkeywords}
Backdoor attacks, large language models, cyber resilience, scaling coefficient,
fine-tuning security, RLHF, adversarial machine learning, MITRE ATLAS.
\end{IEEEkeywords}
